\documentclass[11pt,letterpaper]{article}

% ── Packages ──────────────────────────────────────────────────────────────────
\usepackage[margin=1in]{geometry}
\usepackage{xcolor}
\usepackage[T1]{fontenc}
\usepackage[utf8]{inputenc}
\usepackage[expansion=false,protrusion=false]{microtype}
\usepackage{titlesec}
\usepackage{fancyhdr}
\usepackage{hyperref}
\usepackage{listings}
\usepackage{booktabs}
\usepackage{longtable}
\usepackage{array}
\usepackage{enumitem}
\usepackage{tcolorbox}
\usepackage{parskip}
\usepackage{multirow}
\usepackage{tabularx}

\tcbuselibrary{listings,skins,breakable}

% ── Colors ────────────────────────────────────────────────────────────────────
\definecolor{nfcyan}{HTML}{00FFF0}
\definecolor{nfdark}{HTML}{040C17}
\definecolor{nfyellow}{HTML}{FCEE0A}
\definecolor{nfred}{HTML}{FF2A6D}
\definecolor{nfviolet}{HTML}{7B2FFF}
\definecolor{nfgreen}{HTML}{00FF9F}
\definecolor{codebg}{HTML}{1A1A2E}
\definecolor{codetext}{HTML}{E0E0E0}
\definecolor{cmtgray}{HTML}{888888}
\definecolor{sectionblue}{HTML}{1E3A5F}

% ── Code Listings ─────────────────────────────────────────────────────────────
\lstdefinestyle{shell}{
  backgroundcolor=\color{codebg},
  basicstyle=\ttfamily\small\color{codetext},
  breaklines=true,
  frame=single,
  framerule=0pt,
  rulecolor=\color{nfcyan},
  commentstyle=\color{cmtgray},
  keywordstyle=\color{nfcyan},
  stringstyle=\color{nfgreen},
  showstringspaces=false,
  tabsize=2,
  xleftmargin=6pt,
  xrightmargin=6pt,
}
\lstset{style=shell}

% ── Section Styling ───────────────────────────────────────────────────────────
\titleformat{\section}
  {\large\bfseries\color{white}\colorbox{sectionblue}}
  {\colorbox{sectionblue}{\thesection}}{0.5em}{}
\titleformat{\subsection}{\normalsize\bfseries\color{nfcyan}}{\thesubsection}{0.5em}{}
\titleformat{\subsubsection}{\normalsize\bfseries\color{nfyellow}}{\thesubsubsection}{0.5em}{}

% ── Header/Footer ─────────────────────────────────────────────────────────────
\pagestyle{fancy}
\fancyhf{}
\fancyhead[L]{\small\color{nfcyan}\textbf{NETFRAME} Homelab Runbook}
\fancyhead[R]{\small\color{cmtgray}km-cluster — Kyle Mason}
\fancyfoot[C]{\small\color{cmtgray}\thepage}
\renewcommand{\headrulewidth}{0.4pt}
\renewcommand{\headrule}{\hbox to\headwidth{\color{nfcyan}\leaders\hrule height \headrulewidth\hfill}}

% ── Hyperref ──────────────────────────────────────────────────────────────────
\hypersetup{
  colorlinks=true,
  linkcolor=nfcyan,
  urlcolor=nfcyan,
  pdftitle={NetFRAME Homelab Runbook},
  pdfauthor={Kyle Mason},
}

% ── Info Box ──────────────────────────────────────────────────────────────────
\newtcolorbox{infobox}[1][]{
  colback=codebg, colframe=nfcyan, coltitle=nfdark,
  fonttitle=\bfseries, title=#1, breakable,
}
\newtcolorbox{warnbox}[1][]{
  colback=codebg, colframe=nfred, coltitle=white,
  fonttitle=\bfseries, title=#1, breakable,
}
\newtcolorbox{successbox}[1][]{
  colback=codebg, colframe=nfgreen, coltitle=nfdark,
  fonttitle=\bfseries, title=#1, breakable,
}

% ─────────────────────────────────────────────────────────────────────────────
\begin{document}

% ── Title Page ────────────────────────────────────────────────────────────────
\begin{titlepage}
\pagecolor{nfdark}
\color{white}
\vspace*{2cm}
\begin{center}
  {\Huge\bfseries\color{nfcyan} NETFRAME}\\[0.3cm]
  {\Large\color{white} Homelab Infrastructure Runbook}\\[0.5cm]
  {\large\color{cmtgray} km-cluster \textbar{} NetFRAME CS9000 42U}\\[2cm]
  \rule{0.8\textwidth}{0.4pt}\\[1cm]
  {\normalsize
    \begin{tabular}{rl}
      \color{nfcyan}Author:    & Kyle Mason \\
      \color{nfcyan}GitHub:    & machismo0311/Home-Lab \\
      \color{nfcyan}Revised:   & June 22, 2026 \\
      \color{nfcyan}Cluster:   & km-cluster (7 nodes) \\
      \color{nfcyan}Site:      & kylemason.org \\
    \end{tabular}
  }\\[2cm]
  \rule{0.8\textwidth}{0.4pt}\\[0.5cm]
  {\small\color{cmtgray}
    USMC Aviation Veteran $\cdot$ EC-135/145 Instructor Pilot \\
    Senior FOQA Officer, Metro Aviation \\
    Network Engineering Candidate, VetTec 2.0 / CCNA
  }
\end{center}
\end{titlepage}
\nopagecolor

% ── TOC ───────────────────────────────────────────────────────────────────────
\tableofcontents
\newpage

% =============================================================================
\section{Infrastructure Overview}
% =============================================================================

NetFRAME is a professional-grade homelab built inside a NetFRAME CS9000 42U rack,
purpose-built as a live CCNA lab, LLM inference platform, ML research node,
and backup infrastructure. The cluster is named \textbf{km-cluster} and runs
Proxmox VE 9.1 across seven nodes as of June 2026.

\subsection{Cluster Nodes}

\begin{longtable}{llllll}
\toprule
\textbf{Hostname} & \textbf{Role} & \textbf{IP} & \textbf{CPU} & \textbf{RAM} & \textbf{Notes} \\
\midrule
\endhead
Randy     & Storage / PBS  & 192.168.10.187 & 2x E5-2690 v4 & 128 GB & SuperMicro CSE-219U \\
QuarkyLab & ML / Fernanda  & 192.168.10.30  & 2x E5-2699 v4 & 512 GB & R730, RTX 6000 24GB \\
Jarvis    & LLM Inference  & 192.168.10.31  & 2x E5-2687W v4& 384 GB & R730, RTX 8000 48GB \\
pve2      & OPNsense host  & 192.168.10.204 & i7-8700       & 48 GB  & HP EliteDesk G4 SFF \\
pve3      & Cluster node   & 192.168.10.201 & i7-8700       & 32 GB  & HP EliteDesk G4 SFF \\
pve4      & Cluster node   & 192.168.10.202 & i5-7500T      & 32 GB  & HP EliteDesk G3 Mini \\
pve5      & Cluster node   & 192.168.10.203 & i5-7500T      & 32 GB  & HP EliteDesk G3 Mini \\
\bottomrule
\end{longtable}

\subsection{Management IPs}

\begin{tabular}{ll}
\toprule
\textbf{Device} & \textbf{IP} \\
\midrule
EX3400 (switch) & 192.168.10.50 \\
OPNsense LAN    & 192.168.10.1  \\
Randy IPMI      & 192.168.10.22 \\
QuarkyLab iDRAC & 192.168.10.20 \\
Jarvis iDRAC    & 192.168.10.21 \\
Pi-hole / DNS   & 192.168.10.177 \\
Ares (admin)    & 192.168.10.100 (wired) / .199 (WiFi) \\
\bottomrule
\end{tabular}

% =============================================================================
\section{Network Architecture}
% =============================================================================

\subsection{VLAN Design}

\begin{tabular}{llll}
\toprule
\textbf{VLAN} & \textbf{Name} & \textbf{Subnet} & \textbf{Notes} \\
\midrule
1   & mgmt    & 192.168.10.0/24 & Management, all servers \\
20  & trusted & 192.168.20.0/24 & Trusted devices \\
30  & servers & 192.168.30.0/24 & Service VMs \\
40  & IoT     & 192.168.40.0/24 & IoT devices \\
50  & VoIP    & 192.168.50.0/24 & Cisco 8841 phones \\
60  & guest   & 192.168.60.0/24 & Guest WiFi \\
70  & lab     & 192.168.70.0/24 & CCNA lab segment \\
\bottomrule
\end{tabular}

\subsection{Switching}

\begin{itemize}
  \item \textbf{Juniper EX3400-48P} — enterprise fabric, JunOS 23.4R2-S7.4, IP 192.168.10.50
  \item \textbf{UniFi Switch 24 PRO (PoE+)} — consumer fabric (IoT, VoIP, guest)
  \item Inter-switch uplink: ge-0/0/32 (copper RJ45, active); xe-0/2/3 DAC (link-down, speed mismatch)
\end{itemize}

\subsection{10G Fabric — ConnectX-3 DAC Links}

\begin{tabular}{lll}
\toprule
\textbf{Node} & \textbf{NIC Port} & \textbf{EX3400 Port} \\
\midrule
Randy (nic3)     & Port 1 & xe-0/2/0 — Up 10G \\
Randy (nic2)     & Port 2 & xe-0/2/2 — awaiting 4th DAC \\
QuarkyLab        & ConnectX & xe-0/2/x \\
Jarvis           & ConnectX & xe-0/2/x \\
\bottomrule
\end{tabular}

\subsection{OPNsense}

OPNsense runs as VM 100 on pve2 (192.168.10.204). Version 25.7.
Handles inter-VLAN routing, firewall, DHCP for all VLANs, and WAN.
Serial console access: \texttt{qm terminal 100} from pve2.

% =============================================================================
\section{Randy — SuperMicro CSE-219U}
\label{sec:randy}
% =============================================================================

\begin{infobox}[Randy — Commissioned June 22, 2026]
\textbf{Role:} Primary storage node, Proxmox Backup Server \\
\textbf{IP:} 192.168.10.187 \quad \textbf{IPMI:} 192.168.10.22 \\
\textbf{OS:} Proxmox VE 9.1 (kernel 6.17.2-1-pve) \\
\textbf{CPU:} Dual Intel E5-2690 v4 (2x14c/28t = 48 logical) \\
\textbf{RAM:} 128 GB ECC DDR4
\end{infobox}

\subsection{Storage Layout}

\begin{tabular}{lll}
\toprule
\textbf{Device} & \textbf{Type} & \textbf{Purpose} \\
\midrule
/dev/sda (RAID-1) & 185.8 GB SAS SSD x2 & Proxmox OS boot mirror \\
/dev/sdb--sdu     & 18x 1.6TB Toshiba AL15SEB18EQ 10K SAS & ZFS datastore pool \\
sdc, sdg          & 2x 1.8TB Seagate ST2000NX0423 SATA & Spare / unallocated \\
\bottomrule
\end{tabular}

\subsubsection{RAID Controller — AVAGO 3108 MegaRAID}

\begin{itemize}
  \item Boot mirror: RAID-1 on slots 0:0 and 0:1 (Seagate ST200FM0053 SSDs)
  \item Data drives: JBOD mode enabled, all 20 data drives passed through to OS
  \item StorCLI installed at: \texttt{/opt/MegaRAID/storcli/storcli64}
  \item Symlink: \texttt{/usr/sbin/storcli64}
\end{itemize}

\begin{lstlisting}[language=bash, caption={Useful StorCLI commands}]
# Show all drives and state
storcli64 /c0/eall/sall show

# Show controller info
storcli64 /c0 show

# Enable JBOD mode (if needed after reboot)
storcli64 /c0 set jbod=on
storcli64 /c0/eall/sall set jbod
\end{lstlisting}

\subsubsection{ZFS Pool — datastore}

\begin{lstlisting}[language=bash, caption={ZFS pool layout}]
zpool status datastore
# pool: datastore
# state: ONLINE
#   raidz2-0: sdb sdd sde sdf sdh sdi
#   raidz2-1: sdj sdk sdl sdm sdn sdo
#   raidz2-2: sdp sdq sdr sds sdt sdu
# 29.4TB raw / 19.5TB usable
\end{lstlisting}

\subsection{Proxmox Backup Server}

\begin{itemize}
  \item PBS version: 4.2.2
  \item Web UI: \texttt{https://192.168.10.187:8007}
  \item Datastore: \texttt{datastore} at \texttt{/datastore}
  \item Fingerprint: \texttt{da:61:6a:4c:49:e8:87:03:08:1d:d7:31:ab:23:58:20:47:58:e8:77:4a:52:3d:39:0c:19:52:e0:67:ee:d9:c9}
  \item Storage ID in cluster: \texttt{randy-pbs} (Shared: Yes, all nodes)
\end{itemize}

\subsection{Networking}

\begin{itemize}
  \item \textbf{nic3} (Mellanox ConnectX-3 port 1): 10G UP $\rightarrow$ xe-0/2/0 on EX3400
  \item \textbf{nic2} (ConnectX-3 port 2): DOWN — awaiting 4th DAC cable $\rightarrow$ xe-0/2/2
  \item \textbf{nic0, nic1} (onboard Intel): DOWN / unused
  \item vmbr0 bridges nic3 for management
\end{itemize}

\subsubsection{Randy Commissioning Procedure (for rebuild reference)}

\begin{lstlisting}[language=bash, caption={Full Randy rebuild procedure}]
# 1. Boot IPMI BIOS -- hit Ctrl+H for MegaRAID WebBIOS
#    Create RAID-1 VD on two 185GB Seagate SSDs (slots 0:0, 0:1)

# 2. Install Proxmox VE 9.1 via USB
#    Target: /dev/sda (the RAID-1 VD, 185.78GiB, SMC3108)
#    Filesystem: ext4
#    Hostname: randy.netframe.local
#    IP: 192.168.10.187/24  GW: 192.168.10.1  DNS: 192.168.10.177

# 3. Post-install repo fix
echo "# disabled" > /etc/apt/sources.list.d/pve-enterprise.list
echo "# disabled" > /etc/apt/sources.list.d/ceph.list
echo "" > /etc/apt/sources.list.d/pve-enterprise.sources
echo "" > /etc/apt/sources.list.d/ceph.sources
echo "deb http://download.proxmox.com/debian/pve bookworm pve-no-subscription" \
  > /etc/apt/sources.list.d/pve-no-subscription.list
apt update && apt dist-upgrade -y

# 4. Join cluster
echo "192.168.10.204 pve2.netframe.local pve2" >> /etc/hosts
pvecm add 192.168.10.204 --use_ssh

# 5. Enable JBOD passthrough
storcli64 /c0 set jbod=on
storcli64 /c0/eall/sall set jbod

# 6. Create ZFS pool
zpool create -f -o ashift=12 datastore \
  raidz2 sdb sdd sde sdf sdh sdi \
  raidz2 sdj sdk sdl sdm sdn sdo \
  raidz2 sdp sdq sdr sds sdt sdu

# 7. Install PBS
echo "deb http://download.proxmox.com/debian/pbs trixie pbs-no-subscription" \
  > /etc/apt/sources.list.d/pbs-no-subscription.list
apt update && apt install proxmox-backup-server -y

# 8. Create PBS datastore
proxmox-backup-manager datastore create datastore /datastore
\end{lstlisting}

% =============================================================================
\section{QuarkyLab — Dell R730 \#1}
% =============================================================================

\begin{infobox}[QuarkyLab]
\textbf{Role:} Fernanda's ML research node \\
\textbf{Service Tag:} 1S8WR22 \quad \textbf{iDRAC:} 192.168.10.20 (root/calvin) \\
\textbf{CPU:} Dual E5-2699 v4 (matched stepping) \quad \textbf{RAM:} 512 GB LRDIMM ECC \\
\textbf{GPU:} RTX 6000 24 GB \quad \textbf{BIOS:} 2.19.0
\end{infobox}

\textbf{Pending:} NVIDIA driver install requires kernel pinned to 6.14.11-9-pve before install
(current kernel 6.17 has DRM API incompatibility with nvidia-550 driver).

\begin{lstlisting}[language=bash, caption={Pin kernel before NVIDIA driver install}]
apt install pve-kernel-6.14.11-9-pve
# Set as default in /etc/default/grub, then:
update-grub && reboot
# Then install nvidia-550 driver
\end{lstlisting}

% =============================================================================
\section{Jarvis — Dell R730 \#2}
% =============================================================================

\begin{infobox}[Jarvis]
\textbf{Role:} LLM inference node (Ollama, Qwen2.5 72B) \\
\textbf{Service Tag:} DWG7HH2 \quad \textbf{iDRAC:} 192.168.10.21 \\
\textbf{CPU:} Dual E5-2687W v4 \quad \textbf{RAM:} 384 GB LRDIMM ECC \\
\textbf{GPU:} RTX 8000 48 GB
\end{infobox}

Ollama serves Qwen2.5 72B Q4\_K\_M. FastAPI \texttt{llm\_router.py} implements hybrid
routing — local Ollama by default, escalates to Claude API when logprob confidence
falls below threshold.

% =============================================================================
\section{Services}
% =============================================================================

\subsection{Running Services}

\begin{tabular}{lllll}
\toprule
\textbf{Service} & \textbf{Host} & \textbf{VM/CT ID} & \textbf{IP} & \textbf{Notes} \\
\midrule
OPNsense   & pve2      & VM 100  & 192.168.10.1    & Router/firewall \\
Pi-hole    & pve3      & LXC     & 192.168.10.177  & DNS filter \\
Headscale  & pve3      & LXC 105 & 192.168.10.186  & VPN, v0.29.1 \\
Wazuh      & pve2      & VM 104  & —               & SIEM \\
step-ca    & pve2      & —       & —               & Internal CA \\
PBS        & Randy     & —       & 192.168.10.187:8007 & Backup server \\
\bottomrule
\end{tabular}

\subsection{Headscale}

Headscale v0.29.1 running as LXC 105 on pve3. Ares connected as first node
(Tailscale IP 100.64.0.1). Remaining devices still on commercial Tailscale
pending migration.

\subsection{Internal CA (step-ca)}

step-ca 0.30.2 on pve2 issuing \texttt{*.netframe.local} TLS certificates.

% =============================================================================
\section{Power Infrastructure}
% =============================================================================

\begin{tabular}{lll}
\toprule
\textbf{UPS} & \textbf{Feeds} & \textbf{Battery} \\
\midrule
Middle Atlantic UPS-OL2200R & R730s, Randy, DS4246 & 6x 12V 9Ah AGM in series (76.4V) \\
Tripp Lite SMART1500VA & EX3400, UniFi, small compute & Stock \\
\bottomrule
\end{tabular}

PDU: APC AP7901 (ge-0/0/38 on EX3400, MAC 00:c0:b7:b4:88:4a).

% =============================================================================
\section{Pending / To-Do}
% =============================================================================

\begin{itemize}[leftmargin=*]
  \item[\color{nfred}$\square$] Order 4th DAC cable for Randy nic2 $\rightarrow$ xe-0/2/2
  \item[\color{nfred}$\square$] Pin QuarkyLab kernel to 6.14.11-9-pve, install NVIDIA 550 driver
  \item[\color{nfred}$\square$] Configure backup schedules on all nodes pointing to randy-pbs
  \item[\color{nfred}$\square$] Activate VLANs (pve2 trunk connection to EX3400 pending)
  \item[\color{nfred}$\square$] Migrate Fernanda devices from commercial Tailscale to Headscale
  \item[\color{nfred}$\square$] DS4246 connection to Randy via LSI 9207-8e HBA
  \item[\color{nfred}$\square$] FreePBX VoIP (5x Cisco CP-8841 phones acquired)
  \item[\color{nfred}$\square$] Jellyfin on Randy with RX 580 ROCm transcoding
  \item[\color{nfred}$\square$] UPS NMC IP assignment (requires USB console access)
  \item[\color{nfred}$\square$] Cyberpunk monitoring dashboard live API integration
  \item[\color{nfred}$\square$] RKE2 / Kubernetes planning (Cilium, MetalLB, NVIDIA GPU Operator)
\end{itemize}

% =============================================================================
\section{Quick Reference Commands}
% =============================================================================

\subsection{Cluster Health}

\begin{lstlisting}[language=bash]
# From any node
pvecm status
pvecm nodes

# PBS status
systemctl status proxmox-backup
proxmox-backup-manager datastore list

# ZFS pool health
zpool status datastore
zpool list
\end{lstlisting}

\subsection{EX3400 Switch}

\begin{lstlisting}[language=bash]
ssh mason@192.168.10.50
# Always exit to > operational mode before show commands
# Paste one line at a time

show interfaces xe-0/2/0 brief   # Randy 10G port 1
show interfaces xe-0/2/2 brief   # Randy 10G port 2
show ethernet-switching table     # MAC table
show lldp neighbors               # Connected devices
\end{lstlisting}

\subsection{Randy Storage}

\begin{lstlisting}[language=bash]
ssh root@192.168.10.187

# Check all drives
storcli64 /c0/eall/sall show

# ZFS status
zpool status datastore
zfs list

# If drives go dark after reboot (JBOD mode reset)
storcli64 /c0 set jbod=on
storcli64 /c0/eall/sall set jbod
\end{lstlisting}

\vfill
\begin{center}
\small\color{cmtgray}
NetFRAME Homelab $\cdot$ Kyle Mason $\cdot$ machismo0311/Home-Lab $\cdot$ kylemason.org
\end{center}

\end{document}
